\documentclass[11pt]{article}
\usepackage[a4paper,margin=25mm]{geometry}
\usepackage{amsmath,amssymb,amsthm,longtable,hyperref}
\newtheorem{theorem}{Theorem}\newtheorem{lemma}{Lemma}
\newcommand{\PSL}{\operatorname{PSL}}\newcommand{\Aut}{\operatorname{Aut}}
\newcommand{\Out}{\operatorname{Out}}
\title{A complete small-order domain for the subgroup-atlas reconstruction}
\author{Private L054 recovery component}\date{}
\begin{document}\maketitle
\paragraph{Boundary and retained inputs.}
The fixed 2013 article, Section~4, physical pages~5--6, restricts its
small-order assertion to almost simple groups appearing in its cited
2012 online ATLAS. The final 2015 version, page~264, has the same
qualification. We instead construct a complete mathematical superset.
This does not assert historical file identity. Nor does it establish the
correctness of any unverified, named larger companion table.

We explicitly retain the classified nonabelian-simple-group roster below
$10^6$ documented in GAP~4.15.1, Reference Manual 39.15--15 and
39.15--17, and implemented in the pinned \texttt{grp/simple.gi} and
\texttt{simple.gd}. Its non-$\PSL_2$ data give nineteen entries, of which
$A_5$ and $A_6$ are respectively $\PSL_2(5)$ and $\PSL_2(9)$.
We retain simplicity, the indicated exceptional isomorphisms, and the full
automorphism/outer-group classifications below. For classical groups the
explicit presentations in \texttt{OuterAutoSimplePres} cite
Bray--Holt--Roney-Dougal, pages~36--37; the same pinned implementation
specifies the alternating and sporadic exceptions. These are mathematical
classification inputs, not conclusions inferred from finite order tests.
The ATLAS entries for $L_3(4)$ and $U_3(5)$ independently record the
nonabelian outer groups $D_{12}$ and $S_3$. No new proof of CFSG or of
automorphism classification is claimed.

\begin{lemma}[Preimages and isomorphism classes]
Let $S$ be a nonabelian finite simple group. The almost simple groups with
socle $S$ are exactly the preimages $\pi^{-1}(K)$, where
$\pi:\Aut(S)\to\Out(S)$ and $K\leq\Out(S)$.
Two such preimages are isomorphic if and only if their outer subgroups
are conjugate in $\Out(S)$.
\end{lemma}
\begin{proof}
Identify $S$ with $\operatorname{Inn}(S)$. A group between $S$ and
$\Aut(S)$ containing one representative of an outer coset contains the
whole coset, proving the preimage assertion. Its order is $|S||K|$.
The centralizer of $S$ in $\Aut(S)$ is trivial: if $\alpha$ commutes
with every inner automorphism, then
$c_{\alpha(s)}=\alpha c_s\alpha^{-1}=c_s$, so $\alpha(s)=s$ because
$Z(S)=1$. Any nontrivial normal subgroup $N$ of a preimage meets $S$
nontrivially, for otherwise $[N,S]\leq N\cap S=1$. Simplicity then
implies $S\leq N$. Thus $S$ is the unique minimal normal subgroup
and is characteristic. An isomorphism of two preimages restricts to
$\beta\in\Aut(S)$; conjugation on $S$ shows that its value on any
element $h$ is $\beta h\beta^{-1}$. Therefore the preimages are
conjugate in $\Aut(S)$, exactly when their images are conjugate in
$\Out(S)$. The converse is immediate.
\end{proof}

\begin{lemma}[The linear rank-one part]
There are precisely $86$ almost simple isomorphism classes of order at
most $10^6$ whose socle is $\PSL_2(q)$. They have $39$ distinct simple
socles, taking $q=5$ in place of the duplicate $q=4$.
\end{lemma}
\begin{proof}
For $q=p^f\geq4$ put $d=(2,q-1)$. Counting ordered bases gives
$|\mathrm{GL}_2(q)|=(q^2-1)(q^2-q)$. Determinant is onto
$\mathbb F_q^\times$, and the center of $\mathrm{SL}_2(q)$ has
order $d$. Hence $|S|=q(q^2-1)/d$. The projective action has precisely
the scalar kernel: a linear map fixing every one-dimensional subspace
fixes the two coordinate lines and their sum, forcing equal diagonal
entries. The retained automorphism theorem is
$\Aut(S)=\mathrm{P}\Gamma\mathrm L_2(q)$, including $q=9$.
The determinant modulo squares and the field automorphism give
$\Out(S)=C_d\times C_f$. The field action on the order-two diagonal
quotient is trivial (odd $p$ acts trivially modulo two).

For $q\geq127$, the lower bound $q(q^2-1)/2$ exceeds $10^6$ and
is increasing. Thus it suffices to test prime powers $5\leq q\leq126$.
The literal program \texttt{roster.py} tests primality by trial division,
lists these prime powers, and forms every subgroup of the small explicit
additive group $\mathbb Z/d\times\mathbb Z/f$. Its generic procedure
starts with the identity subgroup and adjoins every possible element to
each known subgroup, taking closure; every finite generating sequence is
therefore represented. A separate check generates every pair of elements.
This also suffices: projection of any subgroup onto $C_f$ is cyclic,
and the kernel is contained in the cyclic group $C_d$; a lift of a
projection generator and a kernel generator generate the subgroup.
The two inventories agree. Keeping precisely those $K$ with
$q(q^2-1)|K|/d\leq10^6$ gives $86$ entries. Since the outer group is
abelian, no further fusion occurs. The only duplicate simple name in
this range is $\PSL_2(4)=\PSL_2(5)$, removed explicitly; the other
listed linear socle orders are distinct.
\end{proof}

\begin{theorem}[Conditional classified roster]
With the explicit classified inputs above, there are exactly $127$
almost simple isomorphism classes of order at most $10^6$. The accepted
eighty current small tables cover eighty of them; the forty-seven
additional targets are all in the $\PSL_2$ family.
\end{theorem}
\begin{proof}
The remaining seventeen simple socles and their retained outer groups
are shown below. $D_{12}$ denotes the dihedral group of order twelve.
The last column counts conjugacy classes of outer subgroups whose
preimages satisfy the bound.
\begin{center}\begin{tabular}{r r c r}
$S$&$|S|$&$\Out(S)$&count\\\hline
$A_7,A_8,A_9$&$2520,20160,181440$&$C_2$&$2+2+2$\\
$L_3(3)$&$5616$&$C_2$&2\\
$L_3(4)$&$20160$&$D_{12}$&10\\
$L_3(5)$&$372000$&$C_2$&2\\
$U_3(3)$&$6048$&$C_2$&2\\
$U_3(4)$&$62400$&$C_4$&3\\
$U_3(5)$&$126000$&$S_3$&4\\
$U_4(2)=\mathrm{PSp}_4(3)$&$25920$&$C_2$&2\\
$\mathrm{Sp}_4(4)$&$979200$&$C_4$&1\\
$\mathrm{Sz}(8)$&$29120$&$C_3$&2\\
$M_{11}$&$7920$&1&1\\
$M_{12}$&$95040$&$C_2$&2\\
$M_{22}$&$443520$&$C_2$&2\\
$J_1$&$175560$&1&1\\
$J_2$&$604800$&$C_2$&1
\end{tabular}\end{center}
The program explicitly multiplies the relevant cyclic and dihedral
groups, closes all subgroups and fuses by every inner conjugation.
Their $41$ surviving classes, together with the $86$ linear classes,
give $127$. The checked roster assigns all eighty accepted targets
once, with exact outer-index multiplicities. Their identities as
the stated almost simple constructions are acquired from their accepted
named-model transports, not deduced from orders. Where several outer
classes have the same order, distinctness is additionally forced by
the accepted numbers of subgroup conjugacy classes:
$25,26,56$ for the three $A_6.2$ groups;
$43,103,39$ for the three $L_2(25).2$ groups;
$151,152,102$ for the three $L_3(4).2$ groups; and
$143,226,142$ for the three index-six $L_3(4)$ extensions.
These are intrinsic isomorphism invariants. Each relevant multiplicity
is therefore exhausted, with no assumption based on a displayed name.
The difference is $127-80=47$.
\end{proof}

\paragraph{Actual finite applications and limits.}
The forty-seven additions comprise twenty prime-field parent descents,
eight nonprime parent descents and nineteen index-two restrictions.
Each uses the previously acquired complete-local-cover descent or
normal-restriction theorem; source counts are never inputs. The projective
models of orders $235200$, $531360$, $885720$, $976500$, $524160$ and $786240$ are literal
defining actions, with full generator and receipt records. Final states
use the accepted explicit generator serializer. The separate two-page applications proof and exact literal certificates
now cover every new row. This roster argument supplies the domain
classification; distinct acceptance of the combined component remains
a separate review step.

Section~4 does not name an individual larger group: it refers to some
larger groups, including large sporadic groups, and then reports an atlas
of more than a hundred outputs. The actual companion index supplies
individual names. A fully verified $127$-table reconstruction would prove
the complete small-order superset and a stronger cardinality statement.
It would not by itself verify every larger PDF named in that index or
prove the exact historical output set. Those claims must retain their
own source-qualified dispositions.
\end{document}
